\makeabstract
    {
    Exploring Hyaluronic Acid-CD44 Signaling in Solid Tumor Cancer Lines  
    }
    {
    Julia Cerimele\textsuperscript{1,2}, Alice Browne\textsuperscript{1}, Miranda Clements\textsuperscript{1}, Wei Ju\textsuperscript{1}, Rosandra Kaplan\textsuperscript{1}
    }
    {
    \textsuperscript{1} Pediatric Oncology Branch, Tumor Microenvironment Section, Center for Cancer Research, National Cancer Institute, National Institutes of Health, Bethesda, MD \\
\textsuperscript{2} Amherst College, Amherst, MA
    }
    {
    The purpose of this project is to investigate whether certain cancer cell lines (F4202, M39M, E0771, Panc02) express CD44, a major receptor of hyaluronic acid (HA). This project further investigates whether different HA molecular weights affect these downstream signaling pathways (Akt, NF-kB, beta-catenin) in cancer cells.
    }
    {
    Data was collected using flow cytometry and western blots. Flow cytometry resuspends a sample in single-cell suspension and passes the cells through a laser, which produces signals indicating cell size and granularity. Fluorescent staining allows for identification of specific targets. For this project, we stained only for CD44. Western blotting separates proteins by size via electrophoresis, transfers them to a membrane, and detects target molecules using antibodies. For the western blots, each cell line was cultured in high (1350 kDa) molecular weight HA, low (40 kDa) molecular weight, or normal media for 24hrs. Image J software was used for pixel intensity analysis of western blots to normalize protein expression.
    }
    {
    Flow cytometry results showed that all cell lines expressed CD44. Western blot analysis confirmed the presence of the correct CD44 isoform. Protein quantification and analysis of the western blot showed that the expression of CD44 was affected by the molecular weight of HA in all cell lines to varying degrees. Additionally, treatment of the Panc02 line with HMW-HA increased NF-kB expression, while treatment with LMW-HA significantly decreased NF-kB expression. Other lines showed varying degrees of NF-kB expression with HMW-HA vs. LMW-HA, but not to the degree seen in the Panc02 line.
    }
    {
    This project showed that different sizes of HA affect CD44 signaling, as well as CD44-related downstream signaling pathways in solid tumor cell lines. Specifically, HMW-HA increases NF-kB expression, and LMW-HA decreases NF-kB in Panc02 lines. Future directions include examining earlier time points in these cell lines to determine when phosphorylation in these pathways occurs, as well as inhibiting these pathways to assess their effects on cancer cell growth/survival for potential future treatments.
    }

    \newpage
    